\section{Factor analysis}
\label{sec:factor_analysis}

\subsection{The orthogonal factor model}
\label{subsec:orthogonal_factor_model}

Let the random vector $\vb{X}$ with $p$ components have mean $\mub$ and
covariance matrix $\mat{S}$ (equivalently $\Si$). The factor model postulates
that the observed variables are linearly dependent on $m$ unobservable random
variables
\[
  F_1,F_2,\dots,F_m \qquad (m<p),
\]
called \textbf{common factors}, and on $p$ additional sources of variation
\[
  \ep_1,\ep_2,\dots,\ep_p,
\]
called \textbf{errors} or \textbf{specific factors}. For $p$ variables one looks
for $m$ factors (with $m<p$) that explain a high percentage of the variation and
that have some meaning in the problem at hand. The model reads

\begin{equation}
  \label{eq:factor_model_rows}
  \small
  \begin{aligned}
    X_1-\mu_1 &= \ell_{11}F_1+\ell_{12}F_2+\dots+\ell_{1m}F_m+\ep_1,\\
    X_2-\mu_2 &= \ell_{21}F_1+\ell_{22}F_2+\dots+\ell_{2m}F_m+\ep_2,\\
              &\;\;\vdots\\
    X_p-\mu_p &= \ell_{p1}F_1+\ell_{p2}F_2+\dots+\ell_{pm}F_m+\ep_p,
  \end{aligned}
\end{equation}

or, in matrix notation,
\begin{equation}
  \label{eq:factor_model_matrix}
  \underset{(p\times 1)}{\vb{X}-\mub}
  =\underset{(p\times m)}{\mat{L}}\;\underset{(m\times 1)}{\vb{F}}
  +\underset{(p\times 1)}{\epsb}.
\end{equation}

The coefficient $\ell_{ij}$ is the \textbf{loading} of the $i$-th variable on
the $j$-th factor, so $\mat{L}$ is the matrix of factor loadings. The $p$
deviations $X_i-\mu_i$ are expressed in terms of $p+m$ random variables
$F_1,\dots,F_m,\ep_1,\dots,\ep_p$, which are unobservable. In summary:

\begin{description}[leftmargin=2.6em,style=nextline,font=\normalfont]
  \item[$\mu_i$] mean of variable $i$.
  \item[$\ep_i$] $i$-th specific factor.
  \item[$F_j$] $j$-th common factor.
  \item[$\ell_{ij}$] loading of the $i$-th variable on the $j$-th factor.
\end{description}

\subsection{Covariance relationships implied by the model}
\label{subsec:factor_covariance_structure}

The orthogonal factor model with $m$ common factors,
\begin{equation}
  \label{eq:factor_model_full}
  \underset{(p\times 1)}{\vb{X}}
  =\underset{(p\times 1)}{\mub}
  +\underset{(p\times m)}{\mat{L}}\;\underset{(m\times 1)}{\vb{F}}
  +\underset{(p\times 1)}{\epsb},
\end{equation}
assumes that $\vb{F}$ and $\epsb$ are independent, with
\begin{align}
  \E(\vb{F})&=\vb{0}, &
  \Cov(\vb{F})&=\E\sqty{\vb{F}\vb{F}^\tp}=\I_m,
  \label{eq:factor_moments_F}\\
  \E(\epsb)&=\vb{0}, &
  \Cov(\epsb)&=\E\sqty{\epsb\epsb^\tp}=\Psib,
  \label{eq:factor_moments_eps}
\end{align}
where $\Psib$ is the diagonal matrix
\begin{equation}
  \label{eq:psi_matrix}
  \Psib=
  \begin{pmatrix}
    \psi_1 & 0 & \cdots & 0\\
    0 & \psi_2 & \cdots & 0\\
    \vdots & \vdots & \ddots & \vdots\\
    0 & 0 & \cdots & \psi_p
  \end{pmatrix},
\end{equation}
and the factors are uncorrelated with the errors,
\begin{equation}
  \label{eq:factor_moments_cross}
  \Cov(\epsb,\vb{F})=\E\sqty{\epsb\vb{F}^\tp}
  =\underset{(p\times m)}{\mat{0}}.
\end{equation}

\selectlanguage{spanish}
\subsection{Planteamiento y consecuencias}
\label{subsec:planteamiento_factorial}

El an\'alisis factorial es una t\'ecnica estad\'istica usada para identificar
patrones ocultos o estructuras latentes detr\'as de grandes conjuntos de
variables observadas. Busca reducir la complejidad agrupando variables
relacionadas en un menor n\'umero de factores subyacentes m\'as f\'aciles de
interpretar.

El contexto es una relaci\'on de la forma
\begin{equation}
  \label{eq:modelo_factorial_es}
  \vb{X}=\mub+\mat{L}\vb{F}+\epsb,
\end{equation}
donde $\vb{X}$ es un vector aleatorio de $p$ variables; $\mub$ un vector
constante de dimensi\'on $p$; $\mat{L}$ una matriz constante $p\times m$, la
``matriz de cargas factoriales''; $\vb{F}\sim\normal(\vb{0},\I_m)$ el vector de
$m$ factores latentes, con media $\vb{0}$ y matriz de covarianzas identidad; y
$\epsb\sim\normal(\vb{0},\Psib)$ el vector de errores \'unicos, con
distribuci\'on normal multivariante de vector de medias $\vb{0}$ y matriz de
covarianzas diagonal $\Psib=\diag\set{\si_1^2,\dots,\si_p^2}$.

Se asume que los factores $\vb{F}$ y los errores $\epsb$ son independientes,
$\Cov(\vb{F},\epsb)=\mat{0}$. Como consecuencia se tiene
\begin{multline}
  \label{eq:esperanza_X_es}
  \E(\vb{X})=\E\qty{\mub+\mat{L}\vb{F}+\epsb}\\
  =\E(\mub)+\mat{L}\,\E(\vb{F})+\E(\epsb)=\mub,
\end{multline}
y para la varianza

\begin{equation}
  \label{eq:varianza_X_es}
  \small
  \begin{aligned}
    \Var(\vb{X})
      &=\Var\qty{\mat{L}\vb{F}+\epsb}\\
      &=\E\sqty{\qty{\mat{L}\vb{F}+\epsb}
                \qty{\mat{L}\vb{F}+\epsb}^\tp}\\
      &=\E\sqty{\mat{L}\vb{F}\vb{F}^\tp\mat{L}^\tp}
        +\E\sqty{\epsb\epsb^\tp}\\
      &\quad+\E\sqty{\mat{L}\vb{F}\epsb^\tp}
        +\E\sqty{\epsb\vb{F}^\tp\mat{L}^\tp}\\
      &=\mat{L}\,\E\sqty{\vb{F}\vb{F}^\tp}\mat{L}^\tp
        +\E\sqty{\epsb\epsb^\tp}\\
      &\quad+\mat{L}\,\E\sqty{\vb{F}\epsb^\tp}
        +\E\sqty{\epsb\vb{F}^\tp}\mat{L}^\tp\\
      &=\mat{L}\I_m\mat{L}^\tp+\Psib
        +\mat{L}\mat{0}+\mat{0}\mat{L}^\tp\\
      &=\mat{L}\mat{L}^\tp+\Psib.
  \end{aligned}
\end{equation}

Como $\vb{X}$ es combinaci\'on lineal de normales, se sigue que
\begin{equation}
  \label{eq:distribucion_X_es}
  \vb{X}\sim\normal\qty{\mub,\;\mat{L}\mat{L}^\tp+\Psib}.
\end{equation}

\subsection{Descomposici\'on de la varianza}
\label{subsec:descomposicion_varianza}

La descomposici\'on $\Var(\vb{X})=\mat{L}\mat{L}^\tp+\Psib$ muestra que la
varianza total de las variables observadas puede repartirse en dos componentes:
\begin{itemize}
  \item \textbf{Varianza com\'un} $\qty{\mat{L}\mat{L}^\tp}$: explicada por los
        factores comunes $\vb{F}$; capta la variabilidad compartida atribuible
        a la estructura latente.
  \item \textbf{Varianza \'unica} $\qty{\Psib}$: espec\'ifica de cada variable
        observada.
\end{itemize}

La varianza de cada variable observada $X_i$ se expresa como
\begin{equation}
  \label{eq:comunalidad}
  \Var(X_i)=\sum_{j=1}^{m}\ell_{ij}^2+\si_i^2,
\end{equation}
donde
\begin{equation}
  \label{eq:h2}
  h_i^2=\sum_{j=1}^{m}\ell_{ij}^2
\end{equation}
es la \textbf{comunalidad} de $X_i$ (la parte explicable por los factores
comunes) y $\si_i^2$ la \textbf{varianza \'unica}.

\subsection{Estimaci\'on de puntuaciones factoriales}
\label{subsec:puntuaciones_factoriales}

Supongamos $\vb{X}=\mub+\mat{L}\vb{F}+\epsb$. La distribuci\'on conjunta de
$\vb{X}$ y $\vb{F}$ es

\begin{equation}
  \label{eq:conjunta_XF}
  \small
  \begin{bmatrix}\vb{X}\\ \vb{F}\end{bmatrix}
  \sim\normal\qty{
    \begin{bmatrix}\mub\\ \vb{0}\end{bmatrix},\;
    \begin{bmatrix}
      \mat{L}\mat{L}^\tp+\Psib & \mat{L}\\
      \mat{L}^\tp & \I_m
    \end{bmatrix}}.
\end{equation}

N\'otese en particular que
\begin{equation}
  \label{eq:cov_XF}
  \begin{aligned}
    \Cov(\vb{X},\vb{F})
      &=\E\sqty{(\vb{X}-\mub)(\vb{F}-\vb{0})^\tp}\\
      &=\E\sqty{\qty{\mat{L}\vb{F}+\epsb}\vb{F}^\tp}\\
      &=\mat{L}\,\E\sqty{\vb{F}\vb{F}^\tp}+\E\sqty{\epsb\vb{F}^\tp}\\
      &=\mat{L}\I_m+\mat{0}=\mat{L}.
  \end{aligned}
\end{equation}

Usando las propiedades de la normal multivariante, se obtiene la distribuci\'on
condicional de $\vb{F}$ dado $\vb{X}$:

\begin{multline}
  \label{eq:condicional_F}
  \sqty{\vb{F}\such\vb{X}=\vb{x}}\sim
  \normal\Big(
    \mat{L}^\tp\qty{\mat{L}\mat{L}^\tp+\Psib}\inv(\vb{x}-\mub),\\
    \I_m-\mat{L}^\tp\qty{\mat{L}\mat{L}^\tp+\Psib}\inv\mat{L}\Big).
\end{multline}

La media de esta distribuci\'on se conoce como la \textbf{puntuaci\'on
factorial} para $\vb{X}=\vb{x}$:
\begin{equation}
  \label{eq:puntuacion_factorial}
  \hat{\vb{F}}
  =\mat{L}^\tp\qty{\mat{L}\mat{L}^\tp+\Psib}\inv(\vb{x}-\mub).
\end{equation}
\selectlanguage{english}
